\documentclass[11pt,a4paper]{article}
\usepackage[margin=25mm]{geometry}
\usepackage{amsmath,amssymb,amsthm,booktabs,hyperref}
\setlength{\emergencystretch}{2em}
\newtheorem{theorem}{Theorem}
\title{Computing the subdivision columns on the original graph}
\author{Human Compression project: a component of the L059 recovery}
\date{}
\begin{document}\maketitle

Let $X$ be a finite connected simple cubic graph, and let $H$ be
an arc-transitive subgroup of its automorphism group. Fix an edge
$e=\{u,v\}$ and put
\[
 P=H_u,\qquad Q=H_e,\qquad B=P\cap Q=H_u\cap H_v.
\]
Here $H_e$ fixes the edge as an unordered set. Permutations act on
the right, as in the accompanying GAP programs, so that cosets of
a subgroup $K$ are written $Kh$.

The subdivision $S(X)$ has original vertices and edge vertices
as its ordered parts, of degrees three and two. Restriction to
the original vertices identifies $\operatorname{Aut}(S(X))$ with
$\operatorname{Aut}(X)$: degrees distinguish the parts, and the
neighbours of each edge vertex determine its two endpoints.
Every automorphism of $X$ extends uniquely. Edges of $S(X)$ are
incident pairs $(u,e)$, or equivalently arcs $(u,v)$ of $X$.
Thus $H$ gives a faithful marking of $S(X)$ with stabilizers
$P,Q,B$. This is the usual graph/coset correspondence used in
De Saedeleer, Leemans, Mixer and Pisanski \cite{source}.

The source's four Max entries describe the actions for
$B<P$, $B<Q$, $P<H$, $Q<H$. An entry is $2T$ for a doubly
transitive action, $M$ for a primitive action that is not doubly
transitive, and $X$ for an imprimitive action.

\begin{theorem}
The following actual permutation actions are respectively
equivalent to the four coset actions defining the Max entries:
\[
 P\curvearrowright X(u),\qquad
 Q\curvearrowright\{u,v\},\qquad
 H\curvearrowright V(X),\qquad
 H\curvearrowright E(X).
\]
Replacing the coset actions by these actions changes none of the
four entries. It changes neither the acting subgroup classification
nor any other invariant column.
\end{theorem}
\begin{proof}
Arc-transitivity of $H$ implies that $P$ is transitive on the three
neighbours of $u$. The stabilizer of $v$ in $P$ is exactly $B$.
Consequently
\[
                   Bp\longmapsto v^p\quad(p\in P)
\]
is a bijection from the right cosets of $B$ in $P$ to $X(u)$.
It is well defined because two representatives give the same
point precisely when their quotient lies in the stabilizer $B$.
It is $P$-equivariant because right multiplication of a coset by
$a$ sends the displayed image to $(v^p)^a=v^{pa}$.

Arc-transitivity also supplies an element of $H$ sending $(u,v)$
to $(v,u)$. It lies in $Q$, so $Q$ is transitive on the two
endpoints of $e$. The stabilizer of $u$ in $Q$ is $B$. The same
argument gives the $Q$-equivariant bijection
$Bq\mapsto u^q$ onto $\{u,v\}$.

Finally, $H$ is transitive on vertices and on edges. Their
stabilizers at $u$ and $e$ are $P$ and $Q$, respectively, so
\[
                  Ph\longmapsto u^h,\qquad
                  Qh\longmapsto e^h
\]
are the other two equivariant bijections. Equivariant bijections
carry blocks to blocks and ordered pairs of distinct points to
ordered pairs of distinct points. They therefore preserve
primitivity and double transitivity, even when an action has a
kernel. This proves equality of all four Max entries.

The new routine receives the same list of actual subgroups $H$,
and computes the same $P,Q,B$. It retains the original normality
tests and greatest element order calculation. A subdivision of a
cubic graph has no type-exchanging automorphism because its two
vertex degrees differ, so its Dual value remains zero. The
subgroup enumeration is executed before this routine and is
unchanged. Thus neither the classified objects nor the remaining
columns are altered.
\end{proof}

\paragraph{Implementation and verification.}
The separate file \texttt{subdivision\_actual\_actions.g} loads the
unchanged original library and replaces only its row routine.
It computes each action image explicitly in GAP, then applies the
same primitivity and double-transitivity tests. It checks the
actual stabilizer intersections, the indices three and two,
generation, faithfulness and transitivity.

The independent comparison script uses three already completed
census witnesses. It recomputes their rows with the new routine
and requires equality of the complete records, including actual
subgroups and all invariant columns:
\begin{center}
\begin{tabular}{lr}
\toprule
Original cubic input & Marking rows compared\\
\midrule
$CAT(1024,10)$ & 3\\
$CAT(1080,1)$ & 4\\
$CAT(1440,6)$ & 3\\
\bottomrule
\end{tabular}
\end{center}
All ten row comparisons pass. This check supplements the uniform
equivariance proof; it is not a new enumeration or a completeness
assertion about those ten rows.

The change was applied only to 319 future worker scripts. The
old worker bytes were retained, and their old and new hashes are
recorded in the transition ledger. Both active runners were on
the same respective input before and after the change. Four
forthcoming inputs for each runner were also left unchanged.
The original library and all completed receipts remain unchanged.
Future receipts record the new worker bytes; the transition ledger
identifies exactly which library those workers load.

The mathematical input is the standard stabilizer/coset
equivariance theorem. Our contribution is its application to
reuse the existing graph actions in this particular computation.
This component makes no new graph classification or whole-original
recovery claim. It uses GAP for the finite permutation computations
and makes no claim of Lean verification.

\begin{thebibliography}{9}
\bibitem{source} J.~De Saedeleer, D.~Leemans, M.~Mixer and
T.~Pisanski, \emph{Core-free, rank two coset geometries from
edge-transitive bipartite graphs}, Mathematica Slovaca
\textbf{64} (2014), 991--1006; arXiv:1106.5704v2.
\end{thebibliography}
\end{document}
